\documentclass[12pt,a4paper]{article}

\usepackage[utf8]{inputenc}
\usepackage[T1]{fontenc}
\usepackage{amsmath, amssymb}
\usepackage{enumitem}
\usepackage{geometry}
\usepackage{xcolor}
\usepackage{multicol}

\geometry{margin=2.3cm}

\title{Banco de Questões Objetivas -- Aula 13\\Técnicas de Controle de Concorrência}
\author{Banco de Dados}
\date{}

\begin{document}

\maketitle

\section*{Instruções}

Este banco contém questões objetivas sobre os principais tópicos da Aula 13:
\begin{itemize}
    \item propósito do controle de concorrência;
    \item conflitos entre transações;
    \item bloqueios;
    \item bloqueio binário;
    \item bloqueio compartilhado e exclusivo;
    \item matriz de compatibilidade;
    \item conversão de bloqueios;
    \item bloqueio em duas fases;
    \item 2PL conservador;
    \item 2PL estrito;
    \item 2PL rigoroso;
    \item deadlock;
    \item starvation.
\end{itemize}

\section{Questões de Múltipla Escolha}

\begin{enumerate}[label=\textbf{Q\arabic*.}]

\item O principal objetivo do controle de concorrência em bancos de dados é:

\begin{enumerate}[label=\alph*)]
    \item permitir que transações conflitantes acessem os mesmos dados sem qualquer regra.
    \item forçar isolamento entre transações conflitantes e preservar a consistência do banco.
    \item eliminar o uso de transações.
    \item substituir o comando \texttt{COMMIT}.
\end{enumerate}

\item O controle de concorrência é necessário porque:

\begin{enumerate}[label=\alph*)]
    \item bancos relacionais não possuem tabelas.
    \item transações concorrentes podem acessar os mesmos itens de dados e gerar inconsistência.
    \item SQL não permite leitura de dados.
    \item toda transação é sempre serial.
\end{enumerate}

\item Um conflito \textit{read-write} ocorre quando:

\begin{enumerate}[label=\alph*)]
    \item duas transações apenas leem o mesmo item.
    \item uma transação lê e outra escreve o mesmo item.
    \item duas transações leem itens diferentes.
    \item uma transação cria uma tabela.
\end{enumerate}

\item Um conflito \textit{write-write} ocorre quando:

\begin{enumerate}[label=\alph*)]
    \item duas transações escrevem no mesmo item de dado.
    \item duas transações apenas leem o mesmo item.
    \item uma transação executa \texttt{SELECT}.
    \item duas transações acessam tabelas diferentes.
\end{enumerate}

\item Diante de conflito entre duas transações sobre um item de dado, o controle de concorrência pode decidir que uma transação:

\begin{enumerate}[label=\alph*)]
    \item sempre será ignorada pelo SGBD.
    \item terá acesso ao item e a outra deverá esperar ou fazer rollback.
    \item deve apagar a tabela inteira.
    \item deve transformar o banco em NoSQL.
\end{enumerate}

\item Um bloqueio é usado para:

\begin{enumerate}[label=\alph*)]
    \item controlar o acesso concorrente a itens de dados.
    \item remover permanentemente uma tabela.
    \item criar uma chave primária.
    \item substituir uma consulta SQL.
\end{enumerate}

\item No bloqueio binário, um item pode estar:

\begin{enumerate}[label=\alph*)]
    \item apenas em modo compartilhado.
    \item bloqueado ou desbloqueado.
    \item apenas em modo exclusivo.
    \item sempre disponível para escrita.
\end{enumerate}

\item Se um item está bloqueado por uma transação em bloqueio binário, outra transação que deseja acessá-lo deve:

\begin{enumerate}[label=\alph*)]
    \item esperar ou ser abortada, dependendo do protocolo.
    \item modificar o item mesmo assim.
    \item apagar o item.
    \item ignorar o bloqueio.
\end{enumerate}

\item Uma limitação do bloqueio binário é que:

\begin{enumerate}[label=\alph*)]
    \item ele permite distinguir bem leitura compartilhada de escrita exclusiva.
    \item ele não diferencia leitura de escrita.
    \item ele elimina deadlocks automaticamente.
    \item ele não usa nenhum estado de bloqueio.
\end{enumerate}

\item Bloqueios compartilhados e exclusivos são usados para:

\begin{enumerate}[label=\alph*)]
    \item permitir maior concorrência entre leituras e controlar escritas.
    \item impedir toda leitura no banco.
    \item remover a necessidade de transações.
    \item substituir dependências funcionais.
\end{enumerate}

\item O bloqueio compartilhado também pode ser chamado de:

\begin{enumerate}[label=\alph*)]
    \item bloqueio de leitura.
    \item bloqueio de escrita.
    \item bloqueio de remoção.
    \item bloqueio de esquema.
\end{enumerate}

\item O bloqueio exclusivo também pode ser chamado de:

\begin{enumerate}[label=\alph*)]
    \item bloqueio de leitura.
    \item bloqueio de escrita.
    \item bloqueio nulo.
    \item bloqueio parcial.
\end{enumerate}

\item Se uma transação deseja apenas ler um item, ela deve solicitar normalmente:

\begin{enumerate}[label=\alph*)]
    \item bloqueio exclusivo.
    \item bloqueio compartilhado.
    \item remoção do bloqueio.
    \item \texttt{ROLLBACK}.
\end{enumerate}

\item Se uma transação deseja escrever em um item, ela deve solicitar normalmente:

\begin{enumerate}[label=\alph*)]
    \item bloqueio exclusivo.
    \item bloqueio compartilhado.
    \item apenas leitura.
    \item \texttt{SELECT DISTINCT}.
\end{enumerate}

\item Dois bloqueios compartilhados sobre o mesmo item são geralmente:

\begin{enumerate}[label=\alph*)]
    \item incompatíveis.
    \item compatíveis.
    \item impossíveis em qualquer SGBD.
    \item equivalentes a rollback.
\end{enumerate}

\item Um bloqueio compartilhado e um bloqueio exclusivo sobre o mesmo item são geralmente:

\begin{enumerate}[label=\alph*)]
    \item compatíveis.
    \item incompatíveis.
    \item sempre permitidos simultaneamente.
    \item irrelevantes para concorrência.
\end{enumerate}

\item Dois bloqueios exclusivos sobre o mesmo item são geralmente:

\begin{enumerate}[label=\alph*)]
    \item compatíveis.
    \item incompatíveis.
    \item obrigatórios para leitura.
    \item permitidos se forem transações diferentes.
\end{enumerate}

\item A matriz de compatibilidade de bloqueios serve para indicar:

\begin{enumerate}[label=\alph*)]
    \item quais bloqueios podem coexistir sobre o mesmo item.
    \item quais tabelas devem ser normalizadas.
    \item quais comandos SQL são DDL.
    \item quais transações já foram apagadas.
\end{enumerate}

\item A conversão de bloqueio de compartilhado para exclusivo é chamada de:

\begin{enumerate}[label=\alph*)]
    \item downgrade.
    \item upgrade.
    \item rollback.
    \item checkpoint.
\end{enumerate}

\item A conversão de bloqueio de exclusivo para compartilhado é chamada de:

\begin{enumerate}[label=\alph*)]
    \item upgrade.
    \item downgrade.
    \item commit.
    \item deadlock.
\end{enumerate}

\item O uso de bloqueios binários, sozinho, garante sempre a serialização dos planos?

\begin{enumerate}[label=\alph*)]
    \item Sim, sempre.
    \item Não, é necessário seguir um protocolo adicional.
    \item Sim, porque todo bloqueio binário é 2PL.
    \item Não, porque bloqueios nunca ajudam na serialização.
\end{enumerate}

\item Para garantir serialização com bloqueios, é necessário controlar principalmente:

\begin{enumerate}[label=\alph*)]
    \item o posicionamento das operações de bloqueio e desbloqueio.
    \item apenas o nome das tabelas.
    \item apenas o número de atributos.
    \item apenas o tipo de dado das colunas.
\end{enumerate}

\item O protocolo de bloqueio em duas fases é conhecido como:

\begin{enumerate}[label=\alph*)]
    \item 1NF.
    \item 2PL.
    \item DDL.
    \item SQLJ.
\end{enumerate}

\item No protocolo de bloqueio em duas fases, todas as operações de \texttt{lock} devem ocorrer:

\begin{enumerate}[label=\alph*)]
    \item depois do primeiro \texttt{unlock}.
    \item antes da primeira operação de \texttt{unlock}.
    \item somente após o \texttt{COMMIT}.
    \item somente depois do \texttt{ROLLBACK}.
\end{enumerate}

\item A primeira fase do bloqueio em duas fases é chamada de:

\begin{enumerate}[label=\alph*)]
    \item fase de encolhimento.
    \item fase de expansão.
    \item fase de recuperação.
    \item fase de normalização.
\end{enumerate}

\item Na fase de expansão de uma transação em 2PL, a transação pode:

\begin{enumerate}[label=\alph*)]
    \item adquirir bloqueios.
    \item liberar bloqueios apenas.
    \item confirmar obrigatoriamente.
    \item apagar o log.
\end{enumerate}

\item A segunda fase do bloqueio em duas fases é chamada de:

\begin{enumerate}[label=\alph*)]
    \item fase de expansão.
    \item fase de encolhimento.
    \item fase de definição.
    \item fase de projeção.
\end{enumerate}

\item Na fase de encolhimento de uma transação em 2PL, a transação pode:

\begin{enumerate}[label=\alph*)]
    \item adquirir novos bloqueios.
    \item liberar bloqueios.
    \item criar novas tabelas.
    \item iniciar outra fase de expansão.
\end{enumerate}

\item Em 2PL, depois que uma transação libera seu primeiro bloqueio:

\begin{enumerate}[label=\alph*)]
    \item ela pode adquirir novos bloqueios livremente.
    \item ela não deve adquirir novos bloqueios.
    \item ela deve apagar todos os dados.
    \item ela deixa de ser uma transação.
\end{enumerate}

\item O protocolo 2PL garante:

\begin{enumerate}[label=\alph*)]
    \item serializabilidade por conflito.
    \item ausência absoluta de deadlock.
    \item ausência absoluta de espera.
    \item impossibilidade de rollback.
\end{enumerate}

\item Uma desvantagem do 2PL básico é que:

\begin{enumerate}[label=\alph*)]
    \item pode ocorrer deadlock.
    \item nunca permite concorrência.
    \item não garante serializabilidade.
    \item elimina todos os bloqueios.
\end{enumerate}

\item Deadlock ocorre quando:

\begin{enumerate}[label=\alph*)]
    \item duas ou mais transações ficam esperando indefinidamente por recursos bloqueados umas pelas outras.
    \item uma transação confirma normalmente.
    \item uma transação executa apenas leitura.
    \item duas transações leem o mesmo item ao mesmo tempo.
\end{enumerate}

\item Um exemplo típico de deadlock é:

\begin{enumerate}[label=\alph*)]
    \item $T_1$ espera por item bloqueado por $T_2$, enquanto $T_2$ espera por item bloqueado por $T_1$.
    \item $T_1$ lê um item e $T_2$ lê o mesmo item.
    \item $T_1$ confirma e $T_2$ inicia.
    \item $T_1$ executa \texttt{SELECT} sem \texttt{WHERE}.
\end{enumerate}

\item Para lidar com deadlock, o SGBD pode:

\begin{enumerate}[label=\alph*)]
    \item abortar uma das transações envolvidas.
    \item remover todas as tabelas do banco.
    \item ignorar todos os bloqueios.
    \item transformar bloqueios exclusivos em valores nulos.
\end{enumerate}

\item Starvation ocorre quando:

\begin{enumerate}[label=\alph*)]
    \item uma transação espera indefinidamente e nunca consegue prosseguir.
    \item uma transação confirma rapidamente.
    \item uma transação não precisa de nenhum bloqueio.
    \item duas transações leem o mesmo item simultaneamente.
\end{enumerate}

\item Starvation pode ocorrer quando:

\begin{enumerate}[label=\alph*)]
    \item uma transação é sempre preterida em favor de outras.
    \item todas as transações são atendidas de forma justa.
    \item nenhuma transação entra em espera.
    \item o banco não possui concorrência.
\end{enumerate}

\item Uma forma de reduzir starvation é:

\begin{enumerate}[label=\alph*)]
    \item usar política justa de escalonamento, como fila.
    \item sempre abortar a mesma transação.
    \item impedir todo bloqueio compartilhado.
    \item apagar o banco periodicamente.
\end{enumerate}

\item O 2PL conservador também é conhecido como:

\begin{enumerate}[label=\alph*)]
    \item 2PL estático.
    \item 2PL nulo.
    \item 2PL sem bloqueios.
    \item 2PL com dirty read.
\end{enumerate}

\item No 2PL conservador, uma transação deve:

\begin{enumerate}[label=\alph*)]
    \item obter todos os bloqueios de que precisa antes de iniciar sua execução.
    \item liberar todos os bloqueios antes de executar qualquer leitura.
    \item nunca solicitar bloqueios.
    \item sempre executar sem esperar.
\end{enumerate}

\item Uma vantagem do 2PL conservador é:

\begin{enumerate}[label=\alph*)]
    \item prevenir deadlocks.
    \item permitir escrita sem bloqueio.
    \item eliminar a necessidade de isolamento.
    \item impedir serializabilidade.
\end{enumerate}

\item Uma desvantagem do 2PL conservador é:

\begin{enumerate}[label=\alph*)]
    \item reduzir a concorrência, pois exige bloqueios antecipados.
    \item permitir deadlock obrigatoriamente.
    \item não usar bloqueios.
    \item impedir toda leitura.
\end{enumerate}

\item No 2PL estrito, os bloqueios exclusivos geralmente são mantidos até:

\begin{enumerate}[label=\alph*)]
    \item o fim da transação.
    \item antes da primeira leitura.
    \item antes de qualquer escrita.
    \item a criação da tabela.
\end{enumerate}

\item O 2PL estrito ajuda principalmente a:

\begin{enumerate}[label=\alph*)]
    \item facilitar a recuperação e evitar cascata de abortos.
    \item eliminar todas as consultas.
    \item remover o uso de \texttt{COMMIT}.
    \item impedir qualquer concorrência de leitura.
\end{enumerate}

\item No 2PL rigoroso, os bloqueios compartilhados e exclusivos são mantidos até:

\begin{enumerate}[label=\alph*)]
    \item o fim da transação.
    \item a primeira operação de leitura.
    \item antes da primeira escrita.
    \item o início da transação.
\end{enumerate}

\item Comparado ao 2PL estrito, o 2PL rigoroso é:

\begin{enumerate}[label=\alph*)]
    \item mais restritivo.
    \item menos restritivo.
    \item idêntico ao bloqueio binário simples.
    \item incompatível com transações.
\end{enumerate}

\item O 2PL estrito ainda pode gerar:

\begin{enumerate}[label=\alph*)]
    \item deadlocks.
    \item ausência total de espera.
    \item ausência de bloqueios.
    \item impossibilidade de serialização.
\end{enumerate}

\item Quando uma transação mantém bloqueios até o \texttt{COMMIT}, isso pode ajudar a evitar:

\begin{enumerate}[label=\alph*)]
    \item leitura de dados não confirmados.
    \item criação de tabelas.
    \item uso de chaves estrangeiras.
    \item consultas com \texttt{GROUP BY}.
\end{enumerate}

\item O controle de concorrência está diretamente relacionado à propriedade ACID de:

\begin{enumerate}[label=\alph*)]
    \item isolamento.
    \item durabilidade apenas.
    \item atomicidade apenas.
    \item normalização.
\end{enumerate}

\item Também contribui para preservar a propriedade de:

\begin{enumerate}[label=\alph*)]
    \item consistência.
    \item 1FN.
    \item BCNF.
    \item criptografia.
\end{enumerate}

\item Em geral, uma transação que não consegue obter um bloqueio compatível deve:

\begin{enumerate}[label=\alph*)]
    \item esperar ou ser abortada, conforme o protocolo.
    \item ignorar a incompatibilidade.
    \item escrever mesmo assim.
    \item apagar o bloqueio de outra transação.
\end{enumerate}

\item Se $T_1$ possui bloqueio compartilhado em $X$, $T_2$ pode obter bloqueio compartilhado em $X$?

\begin{enumerate}[label=\alph*)]
    \item Sim, normalmente é compatível.
    \item Não, nunca.
    \item Sim, mas apenas se for para escrever.
    \item Não, porque leituras sempre conflitam.
\end{enumerate}

\item Se $T_1$ possui bloqueio exclusivo em $X$, $T_2$ pode obter bloqueio compartilhado em $X$ ao mesmo tempo?

\begin{enumerate}[label=\alph*)]
    \item Sim, sempre.
    \item Não, normalmente é incompatível.
    \item Sim, se $T_2$ for mais antiga.
    \item Sim, se $X$ for chave primária.
\end{enumerate}

\item Se $T_1$ possui bloqueio compartilhado em $X$ e deseja escrever em $X$, ela precisa:

\begin{enumerate}[label=\alph*)]
    \item converter o bloqueio compartilhado para exclusivo.
    \item converter o bloqueio exclusivo para compartilhado.
    \item executar \texttt{DROP TABLE}.
    \item remover a transação.
\end{enumerate}

\item Se uma transação começa liberando bloqueios e depois tenta adquirir novos bloqueios, ela viola:

\begin{enumerate}[label=\alph*)]
    \item o protocolo de bloqueio em duas fases.
    \item a 1FN.
    \item a sintaxe de \texttt{SELECT}.
    \item a restrição \texttt{NOT NULL}.
\end{enumerate}

\item Uma execução concorrente correta deve produzir efeito equivalente a:

\begin{enumerate}[label=\alph*)]
    \item alguma execução serial das transações.
    \item uma execução sem transações.
    \item uma execução com todas as transações abortadas.
    \item uma execução sem leituras.
\end{enumerate}

\item Bloqueios ajudam a garantir serializabilidade porque:

\begin{enumerate}[label=\alph*)]
    \item controlam a ordem de acesso a itens conflitantes.
    \item eliminam a necessidade de consistência.
    \item impedem qualquer transação de terminar.
    \item substituem o modelo relacional.
\end{enumerate}

\item A técnica de bloqueio em duas fases está relacionada principalmente a:

\begin{enumerate}[label=\alph*)]
    \item protocolos para garantir propriedades de planos de execução.
    \item criação de atributos multivalorados.
    \item criptografia de senhas.
    \item normalização de relações.
\end{enumerate}

\end{enumerate}

\section{Questões de Verdadeiro ou Falso}

\begin{enumerate}[resume,label=\textbf{Q\arabic*.}]

\item Julgue as afirmações:

\begin{enumerate}[label=\alph*)]
    \item ( \quad ) O controle de concorrência busca preservar a consistência do banco.
    \item ( \quad ) O controle de concorrência ajuda a resolver conflitos entre transações.
    \item ( \quad ) Conflitos \textit{read-write} e \textit{write-write} são relevantes para concorrência.
    \item ( \quad ) Controle de concorrência elimina a necessidade de transações.
\end{enumerate}

\item Julgue as afirmações:

\begin{enumerate}[label=\alph*)]
    \item ( \quad ) Bloqueios controlam acesso a itens de dados.
    \item ( \quad ) Bloqueio binário possui os estados bloqueado e desbloqueado.
    \item ( \quad ) Bloqueio binário diferencia leitura e escrita.
    \item ( \quad ) Bloqueios podem fazer transações esperarem.
\end{enumerate}

\item Julgue as afirmações:

\begin{enumerate}[label=\alph*)]
    \item ( \quad ) Bloqueio compartilhado é usado para leitura.
    \item ( \quad ) Bloqueio exclusivo é usado para escrita.
    \item ( \quad ) Dois bloqueios compartilhados são geralmente compatíveis.
    \item ( \quad ) Dois bloqueios exclusivos sobre o mesmo item são geralmente compatíveis.
\end{enumerate}

\item Julgue as afirmações:

\begin{enumerate}[label=\alph*)]
    \item ( \quad ) Um bloqueio compartilhado é incompatível com bloqueio exclusivo no mesmo item.
    \item ( \quad ) Upgrade converte bloqueio compartilhado em exclusivo.
    \item ( \quad ) Downgrade converte bloqueio exclusivo em compartilhado.
    \item ( \quad ) Conversão de bloqueios nunca ocorre em protocolos de concorrência.
\end{enumerate}

\item Julgue as afirmações:

\begin{enumerate}[label=\alph*)]
    \item ( \quad ) O uso de bloqueio binário sozinho não garante serialização.
    \item ( \quad ) Um protocolo adicional pode ser necessário para garantir serialização.
    \item ( \quad ) O posicionamento de bloqueios e desbloqueios é importante.
    \item ( \quad ) Bloqueios não possuem relação com serializabilidade.
\end{enumerate}

\item Julgue as afirmações:

\begin{enumerate}[label=\alph*)]
    \item ( \quad ) No 2PL, todos os bloqueios devem ser adquiridos antes do primeiro desbloqueio.
    \item ( \quad ) A fase de expansão permite adquirir bloqueios.
    \item ( \quad ) A fase de encolhimento permite liberar bloqueios.
    \item ( \quad ) Depois de liberar um bloqueio, a transação pode adquirir novos bloqueios no 2PL básico.
\end{enumerate}

\item Julgue as afirmações:

\begin{enumerate}[label=\alph*)]
    \item ( \quad ) O 2PL garante serializabilidade por conflito.
    \item ( \quad ) O 2PL básico pode gerar deadlock.
    \item ( \quad ) Deadlock envolve espera circular.
    \item ( \quad ) Deadlock nunca exige intervenção do SGBD.
\end{enumerate}

\item Julgue as afirmações:

\begin{enumerate}[label=\alph*)]
    \item ( \quad ) 2PL conservador tenta obter todos os bloqueios antes da execução.
    \item ( \quad ) 2PL conservador pode prevenir deadlocks.
    \item ( \quad ) 2PL conservador pode reduzir concorrência.
    \item ( \quad ) 2PL conservador dispensa bloqueios.
\end{enumerate}

\item Julgue as afirmações:

\begin{enumerate}[label=\alph*)]
    \item ( \quad ) 2PL estrito mantém bloqueios exclusivos até o fim da transação.
    \item ( \quad ) 2PL estrito ajuda na recuperação.
    \item ( \quad ) 2PL rigoroso mantém bloqueios compartilhados e exclusivos até o fim.
    \item ( \quad ) 2PL rigoroso é menos restritivo que 2PL estrito.
\end{enumerate}

\item Julgue as afirmações:

\begin{enumerate}[label=\alph*)]
    \item ( \quad ) Deadlocks podem ocorrer em protocolos de bloqueio.
    \item ( \quad ) Uma forma de lidar com deadlock é abortar uma transação.
    \item ( \quad ) Starvation ocorre quando uma transação espera indefinidamente.
    \item ( \quad ) Starvation é sempre desejável para melhorar desempenho.
\end{enumerate}

\end{enumerate}

\section{Questões de Aplicação Objetiva}

\begin{enumerate}[resume,label=\textbf{Q\arabic*.}]

\item $T_1$ possui bloqueio compartilhado em $A$. $T_2$ deseja ler $A$. O que normalmente acontece?

\begin{enumerate}[label=\alph*)]
    \item $T_2$ pode obter bloqueio compartilhado em $A$.
    \item $T_2$ deve sempre abortar.
    \item $T_2$ precisa de bloqueio exclusivo.
    \item $T_1$ deve fazer rollback imediatamente.
\end{enumerate}

\item $T_1$ possui bloqueio compartilhado em $A$. $T_2$ deseja escrever em $A$. O que normalmente acontece?

\begin{enumerate}[label=\alph*)]
    \item $T_2$ precisa esperar ou ser abortada, pois precisa de bloqueio exclusivo incompatível.
    \item $T_2$ pode escrever livremente.
    \item $T_2$ só precisa de bloqueio compartilhado.
    \item $T_1$ e $T_2$ podem escrever juntas.
\end{enumerate}

\item $T_1$ possui bloqueio exclusivo em $A$. $T_2$ deseja ler $A$. O que normalmente acontece?

\begin{enumerate}[label=\alph*)]
    \item $T_2$ pode ler normalmente.
    \item $T_2$ deve esperar ou ser abortada.
    \item $T_2$ ganha prioridade automática.
    \item $T_1$ perde o bloqueio sem liberar.
\end{enumerate}

\item Considere a sequência de uma transação:
\[
lock(A),\ read(A),\ unlock(A),\ lock(B),\ read(B),\ unlock(B)
\]
Ela viola 2PL?

\begin{enumerate}[label=\alph*)]
    \item Sim, pois adquire novo bloqueio depois de liberar um bloqueio.
    \item Não, pois toda transação pode alternar lock e unlock livremente.
    \item Não, pois só existem leituras.
    \item Sim, porque não possui escrita.
\end{enumerate}

\item Considere a sequência:
\[
lock(A),\ lock(B),\ read(A),\ write(B),\ unlock(A),\ unlock(B)
\]
Ela respeita 2PL?

\begin{enumerate}[label=\alph*)]
    \item Sim, pois todos os bloqueios foram adquiridos antes do primeiro unlock.
    \item Não, pois 2PL proíbe dois bloqueios.
    \item Não, pois há leitura e escrita.
    \item Sim, mas apenas se não existir \texttt{COMMIT}.
\end{enumerate}

\item $T_1$ bloqueia $A$ e espera por $B$. $T_2$ bloqueia $B$ e espera por $A$. Isso caracteriza:

\begin{enumerate}[label=\alph*)]
    \item deadlock.
    \item starvation apenas.
    \item normalização.
    \item autocommit.
\end{enumerate}

\item Uma transação é continuamente abortada para liberar recursos para outras e nunca consegue terminar. Isso caracteriza:

\begin{enumerate}[label=\alph*)]
    \item starvation.
    \item serialização.
    \item 1FN.
    \item bloqueio compartilhado compatível.
\end{enumerate}

\item Uma transação precisa ler e depois escrever o mesmo item $X$. Inicialmente, ela tem bloqueio compartilhado em $X$. Para escrever, precisa realizar:

\begin{enumerate}[label=\alph*)]
    \item upgrade para bloqueio exclusivo.
    \item downgrade para bloqueio compartilhado.
    \item checkpoint.
    \item \texttt{DROP TABLE}.
\end{enumerate}

\item Qual protocolo seria mais adequado para evitar deadlocks obtendo todos os bloqueios antes de executar?

\begin{enumerate}[label=\alph*)]
    \item 2PL conservador.
    \item 2PL básico sem restrições.
    \item Bloqueio binário sem protocolo.
    \item Execução sem bloqueios.
\end{enumerate}

\item Qual variação mantém bloqueios exclusivos até o final da transação e ajuda a evitar leitura de dados não confirmados?

\begin{enumerate}[label=\alph*)]
    \item 2PL estrito.
    \item 2PL conservador apenas.
    \item Bloqueio binário sem protocolo.
    \item Ausência de bloqueio.
\end{enumerate}

\end{enumerate}

\newpage

\section*{Gabarito}

\begin{multicols}{2}
\begin{enumerate}[label=\textbf{Q\arabic*.}]
    \item b
    \item b
    \item b
    \item a
    \item b
    \item a
    \item b
    \item a
    \item b
    \item a
    \item a
    \item b
    \item b
    \item a
    \item b
    \item b
    \item b
    \item a
    \item b
    \item b
    \item b
    \item a
    \item b
    \item b
    \item b
    \item a
    \item b
    \item b
    \item b
    \item a
    \item a
    \item a
    \item a
    \item a
    \item a
    \item a
    \item a
    \item a
    \item a
    \item a
    \item a
    \item a
    \item a
    \item a
    \item a
    \item a
    \item a
    \item a
    \item a
    \item a
    \item a
    \item a
    \item b
    \item a
    \item a
    \item a
    \item a
    \item a
    \item a
    \item V, V, V, F
    \item V, V, F, V
    \item V, V, V, F
    \item V, V, V, F
    \item V, V, V, F
    \item V, V, V, F
    \item V, V, V, F
    \item V, V, V, F
    \item V, V, V, F
    \item V, V, V, F
    \item a
    \item a
    \item b
    \item a
    \item a
    \item a
    \item a
    \item a
    \item a
    \item a
\end{enumerate}
\end{multicols}

\end{document}
